\documentclass[../main]{subfiles}
\setcounter{chapter}{6}

\begin{document}

\thispagestyle{empty}


\chapter{Alternador trifásico}


\info{Contenido}{Alternadores, frecuencia del alternador, monofásicos, trifásicos, tensiones compuestas, simples, corrientes simples compuestas,carga trifásica}
\actividad{Resolver los siguientes problemas}
\begin{enumerate}
    \item Determinar la frecuencia que produce un alternador que gira a una velocidad de $n=1500 RPM$ si éste posee dos pares de polos.
    \item Determinar la frecuencia que produce un alternador que gira a una velocidad de $n=3500 RPM$ si posee 4 pares de polos.
    \item Determinar el numero de pares de polos que posee un generador que gira a una velocidad de $n=4500rpm$ y genera una frecuencia de 300Hz.
    \item Determinar el numero de pares de polos que posee un generador que gira a una velocidad de $n=3000rpm$ y genera una frecuencia de 100Hz.
    \item ¿A qué velocidad deberá girar un alternador con cuatro pares de polos para producir una frecuencia de $f=50Hz$?
    \item Determinar la tensión compuesta que corresponde a un sistema trifásico que posee una tensión simple de $V_S=223V$.
    \item Determinar la tensión simple que corresponde a un sistema trifásico que posee una tensión compuesta de $V_C=440V$.
    \item Un motor trifásico posee sus bobinas conectadas en estrella. Determinar la corriente eléctrica que absorberá de la línea si al conectarlo a una red con una tensión entre fases de $V_C=400V$ desarrolla una potencia de $P=10KW$ con un factor de potencia de $cos(\varphi)=0,8$ , Averguar la potencia reactiva $Q$ y aparente $S$ del mismo.
    \item Se conectan en estrella tres bobinas iguales a una red trifásica de $V_C=220V, f=50Hz$. Cada una de las mismas posee $R=12\Omega$ de resistencia, y $X_L=20\Omega$ de reactancia inductiva. Calcular la impedancia $Z$, corrientes simples, factor de potencia $cos(\varphi)$, potencia activa $P$ potencia reactiva $Q$ y potencia aparente  $S$.
    \item se desea conectar a una red trifásica con neutro y una tensión entre fases de $V_C=400V$, 30 lámparas fluorescentes de $P=20W,230V cos(\varphi)=0.6$.
    \begin{itemize}
        \item ¿Cómo deben conectarse para que la carga esté equilibrada?
        \item Averiguar la corriente simples y compuestas.
        \item Calcular la potencia del conjunto y por fase.
    \end{itemize}
    \item Un motor trifásico posee sus bobinas conectadas en triángulo. Determinar la corriente eléctrica que absorberá de la línea si al conectarlo a una red con tensión entre fases de $V_C=380V$ , se desarrolla una potencia de $P=15Kw$, con $cos(\varphi)=0,7$. Averiguar la potencia reactiva $Q$ y aparente $S$.
    \item Se conectan en triángulo bobinas de $R=10\Omega$ y $X_L=10\Omega$ a una red trifásica de $V_C=400V, f=50Hz$. Calcular $I_C,I_S, cos(\varphi), P,Q,S$.
    \item Se desean conectar 54 lámparas incandescentes de 100W, 220V, a una red trifásica con una tensión entre fases de $V_C=220V$. ¿Cómo deben conectarse y que consumo poseen?
    \item En un sistema trifásico $V=380V$ con carga equilibrada se mide una intensidad de línea de $I_C=30A$, con un factor de potencia de $cos(\varphi)=0,75$. Si la tensión entre fases es de $V_C=230V$. Averiguar las potencias $Q,P,S$.
    \item En un sistema trifásico con carga equilibrada de tres hilos se mide una potencia en $P=36W$, una intensidad de $I_C=97,4A$. Averiguar el factor de potencia $cos(\varphi)$ de la carga.
    \item Un aparato de calefacción trifásico consta de tres resistencias con $R=10\Omega$ conectadas en estrella. Determinar la potencia que consumirán cuando se aplique $V_C=230V$ entre fases, así como la corriente simple $I_S$ y compuesta $I_C$.
    \item ¿Si conectamos en triángulo las resistencias del problema anterior que valores tomarán $I_C$ y $I_S$?
    \item Un motor trifásico de $P=3990W$ y $cos(\varphi)=0,65$ se conecta a una red de $V_C=380V, f=50hz$, se trata de averiguar la corriente de línea y de cada fase del motor cuando está conectado en triángulo, así como su potencia activa $P$ reactiva $Q$ y aparente $S$. Si cada una de las fases del motor se lo puede considerar como una inductancia $Z$ que posee una reactancia inductiva $X_L$ y una resistencia ohmica $R$. Determina el valor de las mismas.
    \item Se conecta.en triángulo tres bobinas iguales de $R=16\Omega$ y $X_L=10\Omega$ de reactancia inductiva. Si se conectan en un sistema trifásico de $V_C=240V$ entre fases y $f=50Hz$. Determinar
    \begin{itemize}
        \item Las corrientes simples $I_S$ y compuestas $I_C$.
        \item Factor de potencia $cos(\varphi)$.
        \item La potencia Activa $P$ y reactiva $Q$.
    \end{itemize}
    \item Una red trifásica alimenta tres motores monofásicos de inducción de $P=5CV$ ($1CV=736W$), $cos(\varphi)=0,78, V_C=220V$, cada uno conectados entre fase y neutro. Determinar:
    \begin{itemize}
        \item Las corrientes simples y compuestas $I_S,I_C$.
        \item La potencia reactiva que deberá poseer de condensadores  para corregir el $cos(\varphi)=0,9$
    \end{itemize}
    \item Una empresa demanda una potencia aparente $S=700KVA$ a $V_C=10KV$ en corriente alterna trifásica. Las lecturas de consumo en dos meses son para el contador de activa de $Ph=205000KWh$ y $Qh=150000KVAr/h$. Calcular:
    \begin{itemize}
        \item El $cos(\varphi)$ medio en dicho período de facturación.
        \item Intensidades $I_C, I_S$
        \item Potencia reactiva en un banco de capacitores necesaria para corregir el factor de potencia a $cos(\varphi)=0,93$
        \item Porcentaje de reducción de la corriente simple $I_S$ una vez conectado el banco de capacitores.
    \end{itemize}
    \item Una instalación Industrial de 50kW, con un factor de potencia de $cos(\varphi)=0,65$, se alimenta a través de un transformador trifásico, con una tensión en el primario entre fases de $V_{C1}= 24kV$ y de $V_{C2}=400V$ en el secundario. Averiguar:
    \begin{itemize}
        \item La potencia nominal del transformador en kVA, así como la corriente por el primario y secundario.
        \item Nuevas características del transformador si se corrige el a $cos(\varphi)=0.98$ mediante una batería automática de condensadores conectada en el lado de baja tensión.
    \end{itemize}
  \end{enumerate}
\end{document}